% ---------------------------------------------------------------------------
% Abbildung: gemeinsame Test-Schnittstelle IDSPBlock
% benötigt: \usepackage{tikz}
%            \usetikzlibrary{arrows.meta,positioning,calc}
% Einbinden mit: % ---------------------------------------------------------------------------
% Abbildung: gemeinsame Test-Schnittstelle IDSPBlock
% benötigt: \usepackage{tikz}
%            \usetikzlibrary{arrows.meta,positioning,calc}
% Einbinden mit: % ---------------------------------------------------------------------------
% Abbildung: gemeinsame Test-Schnittstelle IDSPBlock
% benötigt: \usepackage{tikz}
%            \usetikzlibrary{arrows.meta,positioning,calc}
% Einbinden mit: % ---------------------------------------------------------------------------
% Abbildung: gemeinsame Test-Schnittstelle IDSPBlock
% benötigt: \usepackage{tikz}
%            \usetikzlibrary{arrows.meta,positioning,calc}
% Einbinden mit: \input{figures/abb_schnittstelle}
% ---------------------------------------------------------------------------
\begin{figure}[htb]
	\centering
	\begin{tikzpicture}[
		font=\footnotesize,
		cls/.style   = {draw, rounded corners=2pt, align=left,
			inner xsep=5pt, inner ysep=4pt, text width=44mm},
		inh/.style   = {draw, -{Triangle[open, length=3mm, width=3mm]}},
		note/.style  = {font=\scriptsize\itshape, align=center}
		]
		
		% ---- abstrakte Schnittstelle -------------------------------------------
		\node[cls, fill=black!14] (iface) at (0,0) {%
			\textbf{\texttt{IDSPBlock}}\\[2pt]
			\texttt{init(sampleRate)}\\
			\texttt{compute(count, in, out)}\\
			\texttt{getNumInputs()}\\
			\texttt{getNumOutputs()}};
		
		% ---- die beiden Implementierungen --------------------------------------
		\node[cls, fill=black!3] (faust) at (-3.4,-3.4) {%
			\textbf{\texttt{FaustBiquad}}\\[2pt]
			hält ein Objekt der vom\\
			Faust-Compiler erzeugten\\
			Klasse \texttt{BiquadDSP} und\\
			reicht die Aufrufe weiter};
		
		\node[cls, fill=black!3] (native) at (3.4,-3.4) {%
			\textbf{\texttt{NativeBiquad}}\\[2pt]
			enthält die manuell\\
			geschriebene Schleife\\
			über die Samples\\
			(Direktform~I)};
		
		\draw[inh] (faust.north)  -- ++(0,0.6) -| (iface.south);
		\draw[inh] (native.north) -- ++(0,0.6) -| (iface.south);
		
		% ---- Nutzer der Schnittstelle ------------------------------------------
		\node[draw, rounded corners=2pt, fill=black!14, align=center,
		inner xsep=5pt, inner ysep=4pt] (user) at (0,2.0)
		{Mess-, Genauigkeits- und Größenprogramme\\
			arbeiten ausschließlich gegen \texttt{IDSPBlock}};
		\draw[-{Stealth[length=2mm]}] (user) -- (iface);
		
	\end{tikzpicture}
	\caption{Gemeinsame Test-Schnittstelle. Beide Varianten erben von derselben
		abstrakten Klasse. Die Messprogramme kennen nur diese Schnittstelle und
		können deshalb für beide Varianten unverändert verwendet werden.}
	\label{fig:schnittstelle}
\end{figure}

% ---------------------------------------------------------------------------
\begin{figure}[htb]
	\centering
	\begin{tikzpicture}[
		font=\footnotesize,
		cls/.style   = {draw, rounded corners=2pt, align=left,
			inner xsep=5pt, inner ysep=4pt, text width=44mm},
		inh/.style   = {draw, -{Triangle[open, length=3mm, width=3mm]}},
		note/.style  = {font=\scriptsize\itshape, align=center}
		]
		
		% ---- abstrakte Schnittstelle -------------------------------------------
		\node[cls, fill=black!14] (iface) at (0,0) {%
			\textbf{\texttt{IDSPBlock}}\\[2pt]
			\texttt{init(sampleRate)}\\
			\texttt{compute(count, in, out)}\\
			\texttt{getNumInputs()}\\
			\texttt{getNumOutputs()}};
		
		% ---- die beiden Implementierungen --------------------------------------
		\node[cls, fill=black!3] (faust) at (-3.4,-3.4) {%
			\textbf{\texttt{FaustBiquad}}\\[2pt]
			hält ein Objekt der vom\\
			Faust-Compiler erzeugten\\
			Klasse \texttt{BiquadDSP} und\\
			reicht die Aufrufe weiter};
		
		\node[cls, fill=black!3] (native) at (3.4,-3.4) {%
			\textbf{\texttt{NativeBiquad}}\\[2pt]
			enthält die manuell\\
			geschriebene Schleife\\
			über die Samples\\
			(Direktform~I)};
		
		\draw[inh] (faust.north)  -- ++(0,0.6) -| (iface.south);
		\draw[inh] (native.north) -- ++(0,0.6) -| (iface.south);
		
		% ---- Nutzer der Schnittstelle ------------------------------------------
		\node[draw, rounded corners=2pt, fill=black!14, align=center,
		inner xsep=5pt, inner ysep=4pt] (user) at (0,2.0)
		{Mess-, Genauigkeits- und Größenprogramme\\
			arbeiten ausschließlich gegen \texttt{IDSPBlock}};
		\draw[-{Stealth[length=2mm]}] (user) -- (iface);
		
	\end{tikzpicture}
	\caption{Gemeinsame Test-Schnittstelle. Beide Varianten erben von derselben
		abstrakten Klasse. Die Messprogramme kennen nur diese Schnittstelle und
		können deshalb für beide Varianten unverändert verwendet werden.}
	\label{fig:schnittstelle}
\end{figure}

% ---------------------------------------------------------------------------
\begin{figure}[htb]
	\centering
	\begin{tikzpicture}[
		font=\footnotesize,
		cls/.style   = {draw, rounded corners=2pt, align=left,
			inner xsep=5pt, inner ysep=4pt, text width=44mm},
		inh/.style   = {draw, -{Triangle[open, length=3mm, width=3mm]}},
		note/.style  = {font=\scriptsize\itshape, align=center}
		]
		
		% ---- abstrakte Schnittstelle -------------------------------------------
		\node[cls, fill=black!14] (iface) at (0,0) {%
			\textbf{\texttt{IDSPBlock}}\\[2pt]
			\texttt{init(sampleRate)}\\
			\texttt{compute(count, in, out)}\\
			\texttt{getNumInputs()}\\
			\texttt{getNumOutputs()}};
		
		% ---- die beiden Implementierungen --------------------------------------
		\node[cls, fill=black!3] (faust) at (-3.4,-3.4) {%
			\textbf{\texttt{FaustBiquad}}\\[2pt]
			hält ein Objekt der vom\\
			Faust-Compiler erzeugten\\
			Klasse \texttt{BiquadDSP} und\\
			reicht die Aufrufe weiter};
		
		\node[cls, fill=black!3] (native) at (3.4,-3.4) {%
			\textbf{\texttt{NativeBiquad}}\\[2pt]
			enthält die manuell\\
			geschriebene Schleife\\
			über die Samples\\
			(Direktform~I)};
		
		\draw[inh] (faust.north)  -- ++(0,0.6) -| (iface.south);
		\draw[inh] (native.north) -- ++(0,0.6) -| (iface.south);
		
		% ---- Nutzer der Schnittstelle ------------------------------------------
		\node[draw, rounded corners=2pt, fill=black!14, align=center,
		inner xsep=5pt, inner ysep=4pt] (user) at (0,2.0)
		{Mess-, Genauigkeits- und Größenprogramme\\
			arbeiten ausschließlich gegen \texttt{IDSPBlock}};
		\draw[-{Stealth[length=2mm]}] (user) -- (iface);
		
	\end{tikzpicture}
	\caption{Gemeinsame Test-Schnittstelle. Beide Varianten erben von derselben
		abstrakten Klasse. Die Messprogramme kennen nur diese Schnittstelle und
		können deshalb für beide Varianten unverändert verwendet werden.}
	\label{fig:schnittstelle}
\end{figure}

% ---------------------------------------------------------------------------
\begin{figure}[htb]
	\centering
	\begin{tikzpicture}[
		font=\footnotesize,
		cls/.style   = {draw, rounded corners=2pt, align=left,
			inner xsep=5pt, inner ysep=4pt, text width=44mm},
		inh/.style   = {draw, -{Triangle[open, length=3mm, width=3mm]}},
		note/.style  = {font=\scriptsize\itshape, align=center}
		]
		
		% ---- abstrakte Schnittstelle -------------------------------------------
		\node[cls, fill=black!14] (iface) at (0,0) {%
			\textbf{\texttt{IDSPBlock}}\\[2pt]
			\texttt{init(sampleRate)}\\
			\texttt{compute(count, in, out)}\\
			\texttt{getNumInputs()}\\
			\texttt{getNumOutputs()}};
		
		% ---- die beiden Implementierungen --------------------------------------
		\node[cls, fill=black!3] (faust) at (-3.4,-3.4) {%
			\textbf{\texttt{FaustBiquad}}\\[2pt]
			hält ein Objekt der vom\\
			Faust-Compiler erzeugten\\
			Klasse \texttt{BiquadDSP} und\\
			reicht die Aufrufe weiter};
		
		\node[cls, fill=black!3] (native) at (3.4,-3.4) {%
			\textbf{\texttt{NativeBiquad}}\\[2pt]
			enthält die manuell\\
			geschriebene Schleife\\
			über die Samples\\
			(Direktform~I)};
		
		\draw[inh] (faust.north)  -- ++(0,0.6) -| (iface.south);
		\draw[inh] (native.north) -- ++(0,0.6) -| (iface.south);
		
		% ---- Nutzer der Schnittstelle ------------------------------------------
		\node[draw, rounded corners=2pt, fill=black!14, align=center,
		inner xsep=5pt, inner ysep=4pt] (user) at (0,2.0)
		{Mess-, Genauigkeits- und Größenprogramme\\
			arbeiten ausschließlich gegen \texttt{IDSPBlock}};
		\draw[-{Stealth[length=2mm]}] (user) -- (iface);
		
	\end{tikzpicture}
	\caption{Gemeinsame Test-Schnittstelle. Beide Varianten erben von derselben
		abstrakten Klasse. Die Messprogramme kennen nur diese Schnittstelle und
		können deshalb für beide Varianten unverändert verwendet werden.}
	\label{fig:schnittstelle}
\end{figure}
